%%%%%%%%%%%%%%%%%%%%%%%%%%%%%%%%%%%%%%%%%
% Beyond Analysis: Strategic Evaluation with Stata
% Eric Booth & Elizabeth Teas
% 2026 Edition
%%%%%%%%%%%%%%%%%%%%%%%%%%%%%%%%%%%%%%%%%

\documentclass[
    a4paper,
    fontsize=10pt,
    twoside=true,
    numbers=noenddot,
]{kaobook}

% --- Language & Typography ---
\ifxetexorluatex
    \usepackage{polyglossia}
    \setmainlanguage{english}
\else
    \usepackage[english]{babel}
\fi
\usepackage[english=british]{csquotes}

% --- Essential Packages ---
\usepackage{blindtext}
\usepackage{kaobiblio}
\IfFileExists{main.bib}{\addbibresource{main.bib}}{} 
\usepackage[framed=true]{kaotheorems}
\usepackage{kaorefs}

% --- Fallback Definitions for Custom kaobook Commands ---
\providecommand{\Environment}[1]{\texttt{#1}}
\providecommand{\Command}[1]{\texttt{\textbackslash#1}}
\providecommand{\Option}[1]{\textsf{#1}}

% Paths for figures
\graphicspath{{images/}} 

% Index & Glossary
\makeindex[columns=3, title=Alphabetical Index, intoc]
\makeglossaries

\begin{document}

%----------------------------------------------------------------------------------------
%	BOOK INFORMATION
%----------------------------------------------------------------------------------------

\titlehead{Tools for Applied Monitoring, Evaluation and Learning}

\title[Beyond Analysis]{Beyond Analysis: \\ Strategic Evaluation with \Environment{Stata}}
\subtitle{subtitle here} 

\author{Eric Booth \hspace{0.5in} Elizabeth Teas} 

\date{Drafted: April 2026}

\publishers{xxxxxxxx Press}

\maketitle

%----------------------------------------------------------------------------------------
%	FRONT MATTER
%----------------------------------------------------------------------------------------

\frontmatter

\makeatletter
\lowertitleback{
    \textbf{Beyond Analysis: Strategic Evaluation with Stata}\\
    © 2025 Eric Booth \& Elizabeth Teas \vspace{0.5in}

    \textbf{Colophon} \\
    This document was typeset using the \texttt{kaobook} class, optimized for the practitioner-researcher. The layout emphasizes marginalia for code snippets and implementation pro-tips.
}
\makeatother

\tableofcontents
\listoffigures
\listoftables

%----------------------------------------------------------------------------------------
%	MAIN BODY
%----------------------------------------------------------------------------------------

\mainmatter
\setchapterstyle{kao}

\chapter*{Preface: The Modern Evaluator's Dilemma}
Applied evaluation has moved beyond simple pre-post t-tests. Today's evaluator must be part data engineer, part AI prompt engineer, and part visualization specialist. This book provides a roadmap for using Stata as the central "brain" of this complex ecosystem.

%----------------------------------------------------------------------------------------
% PART I: THE EMBEDDED EVALUATOR
%----------------------------------------------------------------------------------------

\addpart{Part I: The Embedded Evaluator \newline Principles, Speed, and the Modern Stack}

\chapter{Strategic Evaluation in Implementation Science}
\section{The Practitioner-Researcher Relationship}
Explores the "Embedded Evaluation" model. We discuss moving away from "parachute research" toward sustained relationships where the evaluator is a strategic partner in the implementation team.
\begin{itemize}
    \item \textbf{Bridging the gap}: Shared goals vs. analytical independence. How to maintain rigor while being "in the room" where decisions are made.
    \item \textbf{Feedback loops}: Using \Environment{Stata} to provide real-time insights for program adaptation.
    \item \textbf{Messy Contexts}: Working with nonprofits where data collection is often secondary to service delivery.
\end{itemize}

\section{Implementation Science Principles in Practice}
\marginnote{Implementation science focuses on the 'how' and 'why' of program success, not just the 'what'. See \cite{James2013} for foundational concepts.}
This section defines the core IS frameworks (e.g., RE-AIM, CFIR) and how they translate to data structures in \Environment{Stata}.
\begin{itemize}
    \item \textbf{Fidelity Monitoring}: Tracking if the program is being delivered as intended.
    \item \textbf{Adaptation Tracking}: Coding and managing changes made to the intervention in real-time.
    \item \textbf{Determinants of Change}: Identifying barriers and facilitators through nested data structures.
\end{itemize}

\section{Stata as the Central Evaluation Brain}
Why \Environment{Stata} is the ideal bridge: its ability to handle messy, real-world data while providing the rigor required for publication and the speed required for stakeholders.
\begin{itemize}
    \item \textbf{Reproducibility}: Why do-files are superior to point-and-click dashboards for audit trails.
    \item \textbf{Extensibility}: Using community-contributed commands to solve thorny evaluation problems.
\end{itemize}

\chapter{Setting Up a High-Performance Environment}
\section{The Speed Revolution: \texttt{ftools} and \texttt{gtools}}
\marginnote{When working with millions of administrative records, standard \Environment{Stata} commands can become bottlenecks. We implement \texttt{ftools} and \texttt{gtools} to reclaim your time.}
Detailed guide on using faster alternatives to core commands.
\begin{itemize}
    \item \textbf{Benchmarking}: Comparing \texttt{gcollapse} vs. \texttt{collapse} on large administrative datasets.
    \item \textbf{Multi-way Fixed Effects}: Using \texttt{reghdfe} for complex nested evaluation designs.
    \item \textbf{Efficient Hashing}: Understanding why these tools work and when to use them.
\end{itemize}

\section{Monitoring Real-Time Participation}
Techniques for building "Command Centers" in \Environment{Stata} to track response rates and program reach.
\begin{itemize}
    \item \textbf{Response Rate Trackers}: Automating daily reports from Qualtrics/REDCap APIs.
    \item \textbf{Attrition Alerts}: Using \Environment{Stata} to flag participants at risk of dropping out before they do.
    \item \textbf{Geographic Coverage}: Quick mapping to see where participants are (and aren't) coming from.
\end{itemize}

\section{Communicating with Funders and Stakeholders}
Using \Environment{Stata} to automate the "Monday Morning Update" for nonprofit boards and foundation officers.
\begin{itemize}
    \item \textbf{Dynamic Tables}: Moving beyond copy-pasting to automated \texttt{putexcel} and \texttt{collect} workflows.
    \item \textbf{Funder-Ready Outputs}: Designing visuals that tell a story of impact and reach without the jargon.
\end{itemize}

\section{Integrating AI: The Gemini CLI Bridge}
\marginnote{The \texttt{gemini} package allows \Environment{Stata} to communicate directly with Google's LLMs. We'll cover API key management and \texttt{shell} security.}
Setup guide for integrating Node.js and the Gemini CLI with \Environment{Stata}. 
\begin{itemize}
    \item \textbf{Authentication}: Workflows for macOS, Windows, and Linux.
    \item \textbf{Prompting from Do-files}: Running AI-assisted cleaning routines directly within your pipeline.
    \item \textbf{Privacy}: Strategies for de-identifying data before sending it to an LLM.
\end{itemize}

%----------------------------------------------------------------------------------------
% PART II: DATA ENGINEERING FOR EVALUATORS
%----------------------------------------------------------------------------------------

\addpart{Part II: Data Engineering \newline Modernization, Cleaning, and Harmonization}

\chapter{Advanced Data Collection \& Modernization}
\section{Web Scraping and API Pipelines}
Walks through examples of scraping program data and connecting to APIs like REDCap, Qualtrics, and Census Bureau feeds using \Environment{Stata}.

\section{Unstructured Data: The AI Bridge}
Using LLMs within \Environment{Stata} to turn qualitative interview fragments or messy program notes into structured, analyzable variables using the \texttt{gemini} package.

\chapter{Harmonization and Master Data Management}
\section{Merging the Unmergable}
Fuzzy matching techniques and crosswalk building. Dealing with shifting geographic IDs and changing program taxonomies over time.

\section{Longitudinal Panel Construction}
Building "Long" files from "Wide" chaos. Techniques for tracking program participants across multiple years of intervention in implementation settings.

\chapter{Ethics, Privacy, and Synthetic Data}
\section{The Privacy-Utility Tradeoff}
\marginnote{Data ethics isn't just a compliance box—it's about building trust with participants.}
Techniques for de-identification and the creation of "Research-Ready" public-use files for non-profit and public sector clients.

\section{Synthetic Data Generation}
Using \Environment{Stata}'s \texttt{simulate} and AI-driven synthesis to create realistic mock datasets that preserve the covariance matrix of the original sensitive data.

%----------------------------------------------------------------------------------------
% PART III: THE NEW CAUSAL FRONTIER
%----------------------------------------------------------------------------------------

\addpart{Part III: The New Causal Frontier \newline Innovative and Non-standard Analysis}

\chapter{Modern Difference-in-Differences}
\section{The Revolution in Staggered Adoption}
Explaining the recent shift in DiD methodology and why Two-Way Fixed Effects (TWFE) often fails in applied evaluation. 

\section{Implementing \texttt{csdid} and \texttt{jwdid}}
A step-by-step guide to estimating heterogeneous treatment effects and producing event-study plots that survive academic and client scrutiny.

\chapter{Synthetic Controls and Matrix Completion}
\section{The Synthetic Control Method (\texttt{synth})}
Building a "counterfactual" unit for single-treated unit evaluations. 

\section{Synthetic Difference-in-Differences (\texttt{sdid})}
Combining the strengths of DiD and Synthetic Control for robust program estimation in complex implementation contexts.

\chapter{Machine Learning for Evaluators}
\section{Targeting Models with Lasso and Elastic Net}
Using \texttt{lassopack} to identify which program participants are most likely to drop out or succeed, enabling proactive implementation adjustments.

\section{Random Forests in Stata}
Applying non-linear classification to predict program outcomes and visualize feature importance for stakeholders.

%----------------------------------------------------------------------------------------
% PART IV: HIGH-IMPACT COMMUNICATION
%----------------------------------------------------------------------------------------

\addpart{Part IV: High-Impact Communication \newline Sparklines, Webdocs, and Interactive Reporting}

\chapter{Innovative Visualizations}
\section{Sparklines and Micro-charts with \texttt{sparkta}}
\marginnote{Fahad Mirza's \texttt{sparkta} package allows us to generate thousands of "sparklines" (micro-charts) for high-density dashboards.}
Using \texttt{sparkta} to create dense, data-rich dashboards that show trends at a glance across hundreds of units.

\section{Advanced Coefficient Plots}
Visualizing effect sizes and confidence intervals across multiple models for easy stakeholder comparison.

\chapter{Dynamic Reporting and Web Integration}
\section{Dynamic HTML with \texttt{webdoc}}
\marginnote{By embedding \texttt{sparkta} outputs directly into \texttt{webdoc} workflows, we can generate auto-updating HTML project portals.}
Building live, searchable project websites directly from your do-files to keep practitioner-partners informed.

\section{AI-Assisted Narrative Generation}
Using the \texttt{gemini} package to draft the "Executive Summary" based on your \Environment{Stata} output tables.

%----------------------------------------------------------------------------------------
% PART V: BUILDING TOOLS THAT MATA-ER
%----------------------------------------------------------------------------------------

\addpart{Part V: Building Tools that \texttt{Mata}-er \newline Custom Development and Scalability}

\chapter{Building Internal Evaluation Toolkits}
\section{From Scripts to Systems}
Turning project-specific cleaning code into standardized organization-wide ado-files.

\section{Writing Help Files (\texttt{.sthlp})}
Best practices for documenting institutional knowledge and custom analytical processes for team onboarding.

\chapter{Scalability and Big Data in the Public Sector}
\section{Introduction to Mata for Evaluators}
When even \texttt{gtools} isn't fast enough. Matrix manipulation and custom functions for large-scale administrative data.

%----------------------------------------------------------------------------------------
% BACK MATTER
%----------------------------------------------------------------------------------------

\addpart{Appendices}

\chapter*{Appendix A: The Gemini-Stata Setup Guide}
\chapter*{Appendix B: The Modern Causal Toolbox}

\backmatter
\printbibliography[heading=bibintoc, title=Bibliography]

\end{document}
